\documentclass[exercice,a4]{thedocument}%
\startdatakeys
    \source{laclassecooperativedestef.com}
    \BO{2026}
    \cycle{3}
    \competencesDuSocle{RE4,CO2}
    \objectifsApprentissage{C3C1ab,C3C1aj,C3C2ab,C3C1ae,C3C1ag}
    \motsCles{carre, diagonales, perpendicularite, milieu}
\enddatakeys
\begin{document}%
\noindent\textbf{Programme de construction 4}\par\smallskip\noindent
\begin{itemize}
    \item Tracer le carré $\mathspoint{ABCD}$ de côté $6\text{ cm}$.
    \item Les diagonales se coupent en $\mathspoint{E}$.
    \item Tracer la droite $(\mathspoint{d})$ passant par $\mathspoint{E}$,
    perpendiculaire à $[\mathspoint{CD}]$.
    \item La droite $(\mathspoint{d})$ coupe $[\mathspoint{CD}]$ en
    $\mathspoint{F}$ et $[\mathspoint{BA}]$ en $\mathspoint{G}$.
    \item Placer le point $\mathspoint{H}$ sachant que $\mathspoint{F}$ est
    le milieu de $[\mathspoint{EH}]$.
    \item Placer le point $\mathspoint{I}$ sachant que $\mathspoint{G}$ est
    le milieu de $[\mathspoint{IE}]$.
    \item Tracer les figures $\mathspoint{ECHD}$ et $\mathspoint{IBEA}$,
    puis les colorier en rouge.
    \item Colorier les figures $\mathspoint{BCE}$ et $\mathspoint{AED}$ en
    vert.
    \item Quel est le nom des figures rouges ? Quel est le nom des figures
    vertes ?
\end{itemize}
\end{document}%
